\begin{abstract}
Existing multi-agent environmental decision-support systems rely heavily on black-box Large Language Model (LLM) prompt chaining to arbitrate sensory signals. When modalities suffer from missing values (e.g., satellite cloud occlusion) or conflicting observations, LLM arbiters exhibit uncalibrated confidence, majority-class bias, and a complete lack of source auditability. In this paper, we introduce \textbf{AEGIS} (\textit{Agentic Evidential Geographic Intelligence for Sustainability}), a multi-agent environmental decision-support framework where five domain-specialist agents emit Subjective Logic (SL) opinions over a shared four-state Dirichlet evidence frame $\mathbb{X} = \{\text{Dry}, \text{Saturated}, \text{SurfaceFlow}, \text{Inundation}\}$. A deterministic coordinator performs evidence fusion via Weighted Belief Fusion (WBF) and Consensus \& Compromise Fusion (CCF), dynamically routing predictions based on pairwise Jensen-Shannon conflict divergence and historical Brier-score agent credibility. Missing sensory modalities trigger formal partial-observable projections that decay unobserved inputs into quantified epistemic uncertainty $u$, rather than forcing false confidence. Evaluated on $4,800$ spatio-temporal instances from the Brisbane, Australia 2022 flood event, AEGIS with a 27-dimensional Hybrid Evidential Head achieves \textbf{0.7190 F1-Macro} and \textbf{0.9310 AUROC}, outperforming monolithic late fusion ($0.7106$ F1-Macro) and LLM-arbitrated agentic baselines ($0.2256$ F1-Macro). Furthermore, an analyst-cohort proxy evaluation ($N=12$) confirms $100\%$ explanation completeness and a $+0.82$ Likert trust gain over standard SHAP feature attributions.
\end{abstract}

\begin{IEEEkeywords}
Subjective Logic, Multi-Agent Systems, Evidential Reasoning, Urban Flood Assessment, Uncertainty Quantification, Multimodal Data Fusion.
\end{IEEEkeywords}
